\label{sec:conclusion}

Prison gerrymandering---the Census Bureau's practice of counting
incarcerated people at the prison location rather than their home
address---is a data-layer problem with significant legal and representational
consequences. Approximately 1.4 million people are affected by this
counting convention at each decennial census, and because approximately
58\% of the state and federal prison population is Black or Hispanic,
the distortion falls disproportionately on minority communities. Rural
districts containing large correctional facilities receive inflated
population counts; urban minority communities from which incarcerated
people originate receive deflated counts.

At least twelve states have enacted legislation to correct this distortion
for redistricting purposes, reallocating prison populations to their home
counties using data from the Census Bureau's DHC-A file and the Bureau
of Justice Statistics National Prisoner Statistics program. The
\textsc{bisect} redistricting pipeline implements these corrections through
the \texttt{--adjust-prison-population} flag on \texttt{bisect fetch},
applying the adjustment uniformly across all covered facilities with a
full SHA audit chain.

The empirical findings across North Carolina, Texas, and California
demonstrate that prison population adjustment has measurable but modest
effects on \textsc{bisect} plans: 1--2 districts are affected per state,
with partisan lean shifts of 0.3--0.5 percentage points and minority
voting-age population shifts of 0.4--0.8 percentage points. These shifts
are small relative to the overall plan, but they matter in borderline cases:
a district close to the majority-minority threshold can cross it under
adjusted data, and a district close to the equal-population tolerance
boundary requires a different geographic footprint under the adjusted count.

For expert witnesses and redistricting practitioners, the key takeaway is
procedural: before producing any algorithmic redistricting plan in a state
with a prison population adjustment law, the applicable scope of that law
must be determined (state legislative only, or congressional as well), and
the \textsc{bisect} pipeline must be run with the corresponding data.
The \textit{Evenwel} framework permits this adjustment---it addresses where
within total population to count each person, not which population base
to use---and the \textsc{bisect} implementation provides a reproducible,
judicially verifiable result.

\bigskip
\noindent\textit{The author thanks the redistricting research team for
data access and algorithmic development. All errors are the author's own.}
